% =====================================================================
\section{State of the art}
\label{sec:sota}
% =====================================================================

\paragraph{QKD fundamentals.}
The QKD literature is by now well established. The review by Pirandola
et al.~\cite{pirandola2020} covers discrete-variable protocols (BB84 and
its \emph{decoy-state} variants), continuous-variable protocols and
security proofs. Two constraints that drive all network engineering emerge
from it: the key rate decays exponentially with distance, and a single
link over commercial fibre does not exceed a few hundred kilometres.

\paragraph{From link to network.}
Mehic et al.~\cite{mehic2020} and Cao et al.~\cite{cao2022} describe the
transition from point-to-point links to QKD \emph{networks}: layered
architecture, \emph{trusted relay}, key management and routing. The GÉANT
white paper~\cite{geant2024} gives a readable summary of the same picture
and of the ETSI and ITU-T standards. Routing, which becomes central once a
network is orchestrated, is discussed in~\cite{routing2024}.

\paragraph{Standardisation.}
Two bodies have produced complementary specifications. ETSI ISG QKD started
in 2010 with the use cases~\cite{qkd002} and then defined the interfaces
towards applications~\cite{qkd004,qkd014}, towards the SDN
controller~\cite{qkd015}, towards the orchestrator~\cite{qkd018} and
between key management systems~\cite{qkd020}. ITU-T, with the Y.3800
series, defined the reference architecture of QKD networks
(QKDN)~\cite{ituy3800}, their control and management~\cite{ituy3804} and
SDN-based control~\cite{ituy3805}.

\paragraph{SDN applied to QKD.}
The \emph{Software-Defined Networking} paradigm~\cite{kreutz2015} separates
the control plane from the data plane and centralises decisions in a
programmable controller. The UPM/Telefónica group was the first to
integrate QKD resources into an SDN-controlled optical network, both
experimentally~\cite{aguado2016} and in a systematic
tutorial~\cite{aguado2020}; GS~QKD~015 grew out of this line of work.

\paragraph{Orchestration and multi-domain.}
The orchestrator–controller interface is specified by
GS~QKD~018~\cite{qkd018} and its YANG modules~\cite{qkd018yang}. The first
documented real deployment is that of Sim et al.~\cite{sim2023} (SK
Telecom), who orchestrate two domains from different vendors on the KOREN
testbed. Vicario's thesis~\cite{vicario2025} builds a complete hierarchical
prototype with ONOS and RESTCONF. The MadQCI network~\cite{madqci2024} uses
GS~QKD~015 extensively but describes GS~QKD~018 orchestration as work still
in progress. More recent works propose SDN applications for dynamic link
reconfiguration~\cite{mendez2026}, path-selection criteria and inter-domain
routing that preserves topology confidentiality~\cite{laschet2026}, and
hybrid key-delivery models between regional networks~\cite{interconnect2026}.

\paragraph{Contribution.}
With respect to these sources, this work offers a unified and reasoned
reading of the GS~QKD~002 $\to$ 015 $\to$ 018 chain: it shows how the
requirements of the use cases translate into the components of an SD-QKD
network, explains the structure of the YANG models that describe them and
analyses how GS~QKD~018 abstracts them for multi-domain orchestration.
